\section{Framework and Methods}

\subsection{Architecture Overview}
Spec2Testbench implements a specification-aware verification chain that separates representation, execution, extraction, and decision. The input is a YAML specification parsed into structured entities and then normalized into a finite contract over metrics, units, thresholds, analyses, and optional case-specific metadata. A verification planner transforms this normalized contract into a test intent that binds each requirement to an analysis family, a stimulus pattern, a load condition, one or more observed nodes, and an evidence route. Testbench construction is then performed either deterministically from templates and registry rules or through an optional LLM-assisted formulation stage that proposes testbench text while leaving the evaluation contract unchanged. The generated testbench is executed through a real ngspice backend, after which a result backend selects either native \texttt{.measure} extraction or vector-based \texttt{wrdata} extraction. The resulting values are passed to a metric extractor, then to the specification checker, then to a status classifier that separates execution status, simulation mode, compliance status, robustness status, and scientific category. Finally, a report and provenance layer records the measured values, units, metric traces, backend metadata, and run context in machine-readable and human-readable artifacts.

The implementation follows a Clean Architecture interpretation only in the limited sense required for scientific reproducibility: parsing, planning, generation, simulation, extraction, checking, classification, and reporting are isolated responsibilities with explicit data handoffs. This separation reduces hidden coupling between the circuit under test, the testbench, the simulator, and the decision logic, and it makes it possible to audit whether a conclusion comes from the specification, from the generated testbench, from ngspice execution, or from the post-processing layer.

\subsection{Specification-to-Testbench Translation}
A requirement is translated into an executable testbench contract through a deterministic mapping from semantic intent to simulation evidence. For each expected metric, the planner determines the required analysis type, the excitation or bias stimulus, the load condition, the observed node or signal, the extraction mechanism, the declared unit, and the final assertion applied by the checker. In practice, this means that a specification entry is compiled into an analysis request, one or more source directives, optional load elements, an observation target such as \texttt{vout} or supply current, a native \texttt{.measure} command or a \texttt{wrdata} vector export, a metric name understood by the extractor, a unit normalization path, and a threshold comparison such as range, lower bound, or upper bound.

The translation spans several analysis families. Operating-point and DC requirements are used for quantities such as output bias, quiescent current, or other steady-state voltage and current checks. AC requirements are used for gain- and frequency-related evidence such as \texttt{dc\_gain\_db}, cutoff frequency, unity-gain behavior, or phase-related metrics when these are supported by the specification and backend path. Transient requirements generate time-domain stimuli and timing windows for metrics such as propagation delay, settling behavior, slew-related behavior, and startup amplitude. Spectral requirements use waveform-derived frequency-domain evidence, as in the mixer case where the extractor records quantities such as fundamental frequency or distortion-related outputs. Differential requirements are represented at the specification and planner levels as explicit observed relationships rather than as implicit scalar substitutions. Variation-aware checks are represented as separate requirement families or metadata-scoped scenarios rather than as nominal checks with overloaded semantics; in the present canonical evidence they must be described as a framework capability boundary, not as a completed robustness campaign.

To keep the main text compact, operational command-line usage is intentionally omitted here and can be moved to a compact appendix table. The main methodological point is that the same requirement is preserved across the full path from YAML contract to simulator evidence and back to a checked assertion, so the final decision is tied to an identifiable metric request rather than to a generic simulator success code.

\subsection{Simulation and Evidence Extraction}
All canonical manuscript evidence used here comes from real ngspice execution, not from fabricated or default values. In the nominal paper campaign, all 28 circuits were executed in \texttt{REAL} mode with \texttt{execution\_status = SUCCESS}, and no mock result was retained as scientifically eligible evidence. The execution layer can route extraction through \texttt{NGSPICE\_MEASURE}, which parses native ngspice \texttt{.measure} outputs, or through \texttt{NGSPICE\_WRDATA}, which exports numeric vectors and recomputes supported metrics from the resulting traces. The canonical validation evidence supports both routes: \texttt{NGSPICE\_MEASURE} is exercised in the frozen pilot V3 and the native integration reports, whereas \texttt{NGSPICE\_WRDATA} is validated by two canonical extension cases with independent agreement in both rows of \texttt{results/wrdata\_independent\_comparisons.csv}.

PySpice is optional in this architecture. It is not the authoritative source of the scientific decision, and disabling it does not convert the evaluation into a mock run. Instead, when \texttt{SPEC2TESTBENCH\_DISABLE\_PYSPICE=1}, the framework continues to execute real ngspice runs and uses native extraction and internal parsing paths. The canonical test summary records 66 unique tests passing in both normal mode and PySpice-disabled mode, and the final test artifact explicitly records \texttt{mock\_results\_included = false}. This boundary is important because the paper must distinguish optional convenience tooling from the evidence-producing execution path.

Evidence extraction includes explicit controls against synthetic success. Missing \texttt{.measure} values are kept distinct from zero and propagate as \texttt{NOT\_EVALUATED} rather than as injected numeric defaults. Non-finite values such as NaN or infinity are rejected as non-evaluable measurements. Unit handling is explicit at extraction and checking time so that the final assertion is applied to normalized values with traceable units rather than to raw strings. For oscillatory circuits, frequency acceptance is gated by oscillation validation rather than by the mere presence of an FFT peak. The regression evidence documents this anti-false-PASS behavior directly: missing delay measures now remain null and propagate to \texttt{NOT\_EVALUATED}, and oscillator frequency is blocked unless a validated oscillation criterion is satisfied.

The provenance layer records the conditions under which each measurement was produced. Case-level reports and provenance structures store identifiers for the specification, circuit, testbench, execution mode, status taxonomy, backend choice, measurement source, and run context. This is also where controlled overrides are made auditable: override application status is preserved in metadata and can prevent a result from being silently interpreted as nominal evidence. Scientifically ineligible mock execution, when explicitly permitted for development purposes, is therefore separable from real evaluation in both status and provenance.

\subsection{Deterministic Decision and LLM Boundary}
The final scientific decision in Spec2Testbench is deterministic and is made after ngspice execution by the extraction, checking, and classification pipeline. An optional LLM may assist only at the testbench-construction boundary. Concretely, it may propose an analysis form, a stimulus formulation, parts of the testbench text, or measurement directives that instantiate an already fixed requirement. It may therefore influence how the verification question is expressed to the simulator, but not what counts as compliance.

The LLM is not allowed to modify the benchmark circuit during evaluation, the frozen specification, the thresholds, the backend routing, the parser, the checker, or the final decision rule. It also does not decide whether a missing value should pass, fail, or remain unevaluated. These constraints are essential to the anti-false-PASS discipline. The manuscript should therefore present the LLM, when mentioned at all, as an optional front-end assistant for expressing test intent, while the compliance decision remains anchored in deterministic checking over real simulator evidence with explicit statuses for missing measurements, non-simulable cases, and scientifically ineligible mock outputs.
